\documentclass[11pt]{article}
\usepackage[margin=1in]{geometry}
\usepackage{booktabs}
\usepackage{array}
\usepackage{amsmath}
\usepackage{url}

% This is a static presentation copy generated from docs/cwa_npdes_overview_brief.md,
% which is the canonical, actively-edited source. It is NOT machine-generated by
% a script and has no \input fragments -- if the .md changes, regenerate this
% file by hand (or ask Claude to) rather than editing the two out of sync.

\title{\vspace{-2.5em}The Clean Water Act \& the NPDES Program}
\author{}
\date{\today}

\begin{document}
\sloppy
\maketitle
\vspace{-2.5em}

\section*{Introduction}
This brief situates the National Pollutant Discharge Elimination System (NPDES) within
the broader Clean Water Act (CWA): where the statute came from and how it's organized,
where NPDES (\S\,402) sits among the Act's other major control mechanisms, which states
administer their own NPDES program versus rely on direct EPA permitting, how that
division of authority is visible in the ICIS-NPDES data, and how a violation enters
ICIS.

\section{Origin and structure of the Act}

The CWA's statutory basis dates to the \textbf{Federal Water Pollution Control Act of
1948}, which relied on state-led abatement suits and proved largely ineffective.
Congress rewrote it wholesale in the \textbf{Federal Water Pollution Control Act
Amendments of 1972} (effective October 18, 1972), which introduced technology-based
effluent limits, the NPDES permit program, and the goal of eliminating pollutant
discharges into navigable waters. The 1977 amendments (P.L.\ 95-217) gave the Act its
now-standard short title, ``Clean Water Act.'' The 1987 Water Quality Act added toxics
and stormwater provisions and replaced the Title II municipal construction-grants
program with the Title VI State Revolving Fund.

\paragraph{Purpose.} The Act's objective is to restore and maintain the nation's waters
through making all waters fishable and swimmable by 1983 and having zero water
pollution discharge by 1985. While neither goal was met, these guidelines have greatly
improved waterways nationally and shape the permitting and compliance mechanisms behind
the CWA, specifically NPDES.

The Act is organized into six titles: \textbf{I} (research and related programs),
\textbf{II} (construction grants, now largely superseded by Title VI), \textbf{III}
(standards and enforcement), \textbf{IV} (permits and licenses --- where \S\,402/NPDES
and \S\,404 sit), \textbf{V} (general provisions), and \textbf{VI} (state water
pollution control revolving funds).

\section{The Act's core control mechanisms}

NPDES does not operate in isolation; rather, it fits into one of several interlocking
CWA mechanisms:

\begin{table}[h]
\centering
\begin{tabular}{@{}p{0.08\textwidth}p{0.24\textwidth}p{0.58\textwidth}@{}}
\toprule
Section & Mechanism & What it does \\
\midrule
\S\,301 & Effluent limitations & Bans discharging without a permit; sets
  technology-based limits \\
\S\,303 & Water quality standards / TMDLs & States set standards per water body;
  impaired waters get a TMDL \\
\S\,304 & Criteria \& guidelines & EPA's technical criteria and effluent guidelines
  feed into permit limits \\
\S\,319 & Nonpoint source management & Grants only, no permit --- why nonpoint runoff
  is outside NPDES \\
\S\,401 & State certification & State sign-off that a federal permit (\S\,402 or
  \S\,404) meets state water quality standards \\
\S\,402 & NPDES & Point-source discharge permits --- this project's focus \\
\S\,404 & Dredge and fill & Army Corps' permit for dredged/fill material, separate
  from \S\,402 \\
\bottomrule
\end{tabular}
\caption{Core CWA control mechanisms and how they relate to NPDES.}
\end{table}

A \S\,402 permit's limits combine \S\,301's technology-based floor with
\S\,303/\S\,304's water-quality-based ceiling --- the basis for the
``technology-based \textit{and} water-quality-based limits'' language in
\path{permit_types_brief.md} \S\,2. And \S\,401 certification isn't NPDES-specific: it
applies to \S\,404 permits just as much.

\section{NPDES's place in this structure}

Within \S\,402, NPDES is the mechanism that authorizes a point-source discharge ---
everything else in the Act (\S\,301's baseline, \S\,303/\S\,304's water-quality inputs,
\S\,401's state sign-off) determines what a given NPDES permit may require.

\section{Who administers NPDES: state primacy vs.\ direct federal implementation}
\label{sec:delegation}

CWA \S\,402(b) lets a state apply to administer its own NPDES program in lieu of EPA,
subject to EPA approval and continuing oversight (EPA retains authority to object to
individual state-issued permits and to step in on enforcement). Once approved, the
state issues and administers permits directly, though records for both state-issued
and EPA-issued permits are tracked in the same shared ICIS-NPDES system.

As of 2026, 47 states are fully authorized to run their own NPDES program, while
Massachusetts, New Hampshire, and New Mexico are run by the EPA. Permits for
Washington, D.C.\ and nearly every U.S.\ territory (the Virgin Islands is the one
territory that runs its own program) are also federally controlled. Indian country is
a separate case: EPA retains permitting authority there in every state by default,
regardless of whether that state is authorized, though a tribe can separately apply
for its own authority under \S\,518.

\section{How authorization status appears in the data}

Delegation status is \textbf{not} a first-class field on most ICIS-NPDES permit or
facility records --- a permit's \texttt{STATE\_CODE} tells you where a facility is, not
who issued the permit, and there is no simple ``issuing authority'' flag on
\texttt{ICIS\_PERMITS.csv} itself. It does surface in two narrower places:

\begin{itemize}
  \item \textbf{Enforcement action type codes}: several \texttt{ENF\_TYPE\_CODE}
        values come in state/EPA pairs distinguished only by a trailing \texttt{S}
        (e.g.\ \texttt{AER}/\texttt{AERS}, \texttt{PHEMAIL}/\texttt{PHEMLS}) --- a
        record-level flag for \textit{who took this specific enforcement action},
        separate from the permit's own issuing authority.
  \item \textbf{The Master General Permits download}: each \texttt{NGP} master record
        carries an explicit \texttt{ISSUING\_AGENCY} field, since a master general
        permit is issued once by a single agency (state or EPA) rather than inherited
        implicitly from wherever a covered facility happens to sit.
\end{itemize}

For everything else --- in particular, individual (\texttt{NPD}) permits, which is
where most project analysis concentrates --- whether a given permit was issued by a
state or by EPA has to be inferred from the state list in Section~\ref{sec:delegation}
above, not read off a field in the data.

\section{How a violation enters ICIS}

A violation reaches ICIS through one of two detection pathways, and the data doesn't
always make clear which:

\begin{itemize}
  \item \textbf{Self-reported noncompliance} --- the facility's own Discharge
        Monitoring Report shows it exceeded a permit limit. This is the routine channel
        and the source of the large majority of effluent-violation rows.
  \item \textbf{Independent inspection} --- a sampling inspection catches an exceedance
        the facility did not itself report. This channel is comparatively rare,
        concentrated among individual/major permits, and tied to the routine-inspection
        and SNC apparatus described in \path{permit_types_brief.md} \S\,6.
\end{itemize}

\section{Implications for analysis}

Two caveats follow directly from Sections~4--6.

\begin{itemize}
  \item \textbf{Enforcement and violation counts are not directly comparable across
        the 47 delegated states and the EPA-direct jurisdictions} (Massachusetts, New
        Hampshire, New Mexico, the District of Columbia, most territories, and Indian
        country nationwide). State agencies and EPA regional offices differ
        systematically in enforcement discretion, staffing levels, and institutional
        priorities; an observed difference between, for instance, Massachusetts and a
        neighboring authorized state may reflect this administrative distinction as
        much as underlying facility behavior. Because jurisdiction type is not itself a
        field in the permit data, this distinction must be supplied exogenously, via
        the classification in Section~\ref{sec:delegation}.
  \item \textbf{Violation counts are a function of monitoring and inspection
        intensity, not solely of noncompliance.} Given that self-reported DMR
        exceedances substantially outnumber inspection-detected ones (Section~6), a
        facility, state, or permit type subject to less frequent monitoring will, by
        construction, register fewer violations irrespective of its true compliance
        rate.
\end{itemize}

\section*{Reproducibility}
This document is a static presentation copy of \path{docs/cwa_npdes_overview_brief.md}
(the actively-edited, canonical source --- edit that file, not this one). This brief
draws statutory and regulatory content from 33 U.S.C.\ Chapter 26 (CWA \S\S\,101, 301,
303, 304, 319, 401, 402, 404, 518), 40 CFR Part 123, EPA's ``NPDES State Program
Authority'' program page, the Federal Register notice approving Idaho's NPDES program
(83 Fed.\ Reg.\ 27668, June 14, 2018) and its associated phased-implementation
schedule, and reporting on Massachusetts's delegation efforts (Foley Hoag, ``And Then
There Were Three,'' 2018; Massachusetts Municipal Association delegation updates).

\end{document}
